% =========================================================================
% UNIT 9: INFERENCE FOR QUANTITATIVE DATA - SLOPES
% =========================================================================
\chapter{Unit 9: Inference for Quantitative Data: Slopes}

\section{Unit Overview \& AP Exam Weight}
Unit 9 accounts for \textbf{2\% to 5\%} of the multiple-choice section of the AP Statistics Exam and appears frequently on Section II Free-Response Questions (often as part of Question 1 or Question 6 Investigative Task). It integrates two-variable regression analysis from Unit 2 with formal inference concepts ($t$-intervals and $t$-tests).

\subsection*{Key Vocabulary Terms}
\begin{description}[leftmargin=1.5em, style=nextline]
  \item[Population Regression Line:] $\mu_y = \alpha + \beta x$. Models the true mean response $\mu_y$ for any given value of explanatory variable $x$. Parameter $\beta$ represents the true population slope.
  \item[Sample Regression Line (LSRL):] $\hat{y} = b_0 + b_1 x$. The estimated regression line calculated from sample data, where $b_1$ is the point estimate of $\beta$ and $b_0$ estimates $\alpha$.
  \item[Standard Error of the Slope ($\text{SE}(b_1)$):] The estimated standard deviation of the sampling distribution of sample slopes $b_1$:
  \begin{equation}
    \text{SE}(b_1) = \frac{s}{s_x \sqrt{n - 1}}
  \end{equation}
  where $s$ is the standard deviation of the residuals and $s_x$ is the sample standard deviation of $x$.
  \item[Standard Deviation of Residuals ($s$):] Estimates the standard deviation $\sigma$ of the response variable $y$ about the true population regression line:
  \begin{equation}
    s = \sqrt{\frac{\sum (y_i - \hat{y}_i)^2}{n - 2}} = \sqrt{\frac{\sum e_i^2}{n - 2}}
  \end{equation}
  \item[Degrees of Freedom for Regression:] $df = n - 2$, reflecting the estimation of two parameters ($\alpha$ and $\beta$).
  \item[LINER Conditions:] The five mandatory assumptions for valid regression inference: Linear, Independent, Normal residuals, Equal variance of residuals, Random sample/assignment.
\end{description}

\section{Core Concept Lessons}

\begin{lessonbox}[Lesson 9.1: The L-I-N-E-R Conditions for Regression Inference]
On an AP exam FRQ, you must state and verify all 5 LINER conditions:
\begin{enumerate}[leftmargin=1.5em]
  \item \textbf{L --- Linear:} The true relationship between $x$ and $y$ in the population is linear.\\
  \textit{Verification:} The scatterplot shows a roughly linear pattern, and the residual plot exhibits random scatter around the horizontal zero line with no curved pattern.
  \item \textbf{I --- Independent:} Individual observations are independent.\\
  \textit{Verification:} When sampling without replacement, verify $n \le 0.10N$. In an experiment, verify subjects were randomly assigned to treatments.
  \item \textbf{N --- Normal Residuals:} For any fixed $x$, responses $y$ vary normally about the line.\\
  \textit{Verification:} Construct a histogram, boxplot, or normal probability plot of the residuals. Confirm there is no extreme skewness and no severe outliers.
  \item \textbf{E --- Equal Variance (Homoscedasticity):} The standard deviation of $y$ ($\sigma$) is constant across all values of $x$.\\
  \textit{Verification:} In the residual plot, the vertical spread of the residuals is roughly uniform across the range of $x$ values (no fan-shaped or funnel pattern).
  \item \textbf{R --- Random:} Data were collected using a well-designed random sample or randomized experiment.\\
  \textit{Verification:} Quote the problem stem confirming an SRS or randomized experimental protocol.
\end{enumerate}
\end{lessonbox}

\begin{lessonbox}[Lesson 9.2: Confidence Interval for Slope $\beta$]
A $C\%$ confidence interval for the true population slope $\beta$ is:
\begin{equation}
  b_1 \pm t^* \cdot \text{SE}(b_1)
\end{equation}
where $t^*$ is the critical value from the $t$-distribution with $df = n - 2$.

\begin{frqrubricbox}[Score-4 Template: Interpreting a Confidence Interval for Slope]
\textit{"We are [C\%] confident that for each additional 1 [unit of $x$] increase in [explanatory variable $x$], the true population mean [response variable $y$] increases/decreases by an amount between [Lower Bound] and [Upper Bound] [units of $y$]."}
\end{frqrubricbox}
\end{lessonbox}

\begin{lessonbox}[Lesson 9.3: Hypothesis Test for Slope $\beta$ ($H_0: \beta = 0$)]
To test whether there is a statistically significant linear association:
\begin{align*}
  H_0 &: \beta = 0 \quad (\text{There is no linear association between } x \text{ and } y \text{ in the population}) \\
  H_a &: \beta \neq 0 \quad \text{or} \quad \beta > 0 \quad \text{or} \quad \beta < 0
\end{align*}
The test statistic is:
\begin{equation}
  t = \frac{b_1 - 0}{\text{SE}(b_1)}, \qquad df = n - 2
\end{equation}
If $p\text{-value} \le \alpha$, reject $H_0$ and conclude that convincing evidence exists of a linear relationship between $x$ and $y$.
\end{lessonbox}

\begin{lessonbox}[Lesson 9.4: Deciphering Computer Regression Output]
AP exams regularly provide regression computer output rather than raw data. You must extract key statistics instantly:
\begin{center}
\begin{adjustbox}{max width=\linewidth}
\begin{tabular}{lcccc}
\toprule
\textbf{Predictor} & \textbf{Coef} & \textbf{SE Coef} & \textbf{T} & \textbf{P} \\
\midrule
Constant & $b_0$ ($y$-intercept) & $\text{SE}(b_0)$ & $t = b_0 / \text{SE}(b_0)$ & $p\text{-value}$ \\
\textbf{Explanatory Variable ($x$)} & \textbf{$b_1$ (slope)} & \textbf{$\text{SE}(b_1)$} & \textbf{$t = b_1 / \text{SE}(b_1)$} & \textbf{Two-sided $p$-value} \\
\bottomrule
\end{tabular}
\end{adjustbox}
\end{center}
\textbf{Additional Output Statistics:}
\begin{itemize}
  \item $S = s$: The standard deviation of the residuals. Measures the typical prediction error.
  \item $R\text{-sq} = r^2$: The coefficient of determination.
  \item \textbf{Direction of $r$:} The computer output reports $r^2$ as a positive percentage. To find $r$, calculate $\pm \sqrt{r^2}$, giving $r$ the same sign as the slope coefficient $b_1$!
  \item \textbf{One-Sided $p$-Value Warning:} Computer tables report \textbf{two-sided} $p$-values. If your alternative hypothesis is one-sided ($H_a: \beta > 0$ or $H_a: \beta < 0$), divide the table $p$-value by 2 (provided the sample slope is in the hypothesized direction)!
\end{itemize}
\end{lessonbox}

\section{Worked Step-by-Step Examples}

\subsection*{Example 9.1: Regression Output Analysis \& Interpretation}
\textbf{Problem:} An automotive research team investigates the relationship between vehicle curb weight ($x$, in thousands of pounds) and highway fuel efficiency ($y$, in miles per gallon). An SRS of $n = 18$ passenger cars yields:
\begin{center}
\begin{tabular}{lcccc}
\toprule
\textbf{Predictor} & \textbf{Coef} & \textbf{SE Coef} & \textbf{T} & \textbf{P} \\
\midrule
Constant & 48.24 & 2.18 & 22.13 & 0.000 \\
Weight & $-5.12$ & 0.64 & $-8.00$ & 0.000 \\
\bottomrule
\end{tabular}
\end{center}
$S = 1.86\text{ mpg}, \quad R\text{-sq} = 80.0\%$.\\
(a) Write the equation of the least-squares regression line.\\
(b) Interpret the slope $b_1 = -5.12$ in context.\\
(c) Interpret $S = 1.86\text{ mpg}$ in context.\\
(d) Construct a 95\% confidence interval for the true slope $\beta$.

\textbf{Solution:}
\begin{itemize}[leftmargin=1.2em]
  \item \textbf{Part (a):} $\widehat{\text{Fuel Efficiency}} = 48.24 - 5.12(\text{Weight})$.
  \item \textbf{Part (b):} For each additional 1,000-pound increase in vehicle weight, the predicted highway fuel efficiency decreases by approximately 5.12 miles per gallon.
  \item \textbf{Part (c):} When using vehicle weight to predict highway fuel efficiency with this model, predictions will typically be off by about 1.86 miles per gallon from the actual observed fuel efficiency.
  \item \textbf{Part (d):} Degrees of freedom: $df = n - 2 = 18 - 2 = 16$.
  From the $t$-table, for 95\% confidence with $df = 16$, $t^* = 2.120$.\\
  \[ \text{CI} = b_1 \pm t^* \cdot \text{SE}(b_1) = -5.12 \pm 2.120(0.64) = -5.12 \pm 1.3568 = \mathbf{(-6.477, \; -3.763)\text{ mpg/1000 lbs}} \]
  \textit{"We are 95\% confident that for each additional 1,000-pound increase in vehicle curb weight, the true mean highway fuel efficiency decreases by between 3.76 and 6.48 miles per gallon."}
\end{itemize}

\subsection*{Example 9.2: Hypothesis Test for Positive Linear Association}
\textbf{Problem:} Using the tree dataset from Example 9.1 where $n = 16$, $b_1 = 2.85$, and $\text{SE}(b_1) = 0.42$, carry out a formal hypothesis test to determine if there is convincing evidence that taller trees have larger diameters ($\beta > 0$) at $\alpha = 0.01$.

\textbf{Solution:}
\begin{itemize}[leftmargin=1.2em]
  \item \textbf{Step 1 (Hypotheses):}
  $H_0: \beta = 0$ (No positive linear relationship between diameter and height).\\
  $H_a: \beta > 0$ (There is a positive linear relationship between diameter and height).
  \item \textbf{Step 2 (Conditions):}
  Assume all LINER conditions are met as verified by scatterplots and residual plots.
  \item \textbf{Step 3 (Mechanics):}
  \[ t = \frac{b_1 - 0}{\text{SE}(b_1)} = \frac{2.85 - 0}{0.42} = 6.786 \approx 6.79 \]
  Degrees of freedom: $df = 16 - 2 = 14$.\\
  $p$-value: $P(t_{14} \ge 6.79) = \text{tcdf}(6.79, 10^{99}, 14) \approx 0.000004 < 0.0001$.
  \item \textbf{Step 4 (Conclusion):}
  Because the $p$-value ($< 0.0001$) is much less than $\alpha = 0.01$, we reject $H_0$. There is convincing statistical evidence of a positive linear association between pine tree diameter and tree height.
\end{itemize}

\section{TI-84 CE Calculator Procedures for Slope Inference}
\begin{calculatorbox}[TI-84 Playbook: Linear Regression Inference]
\begin{itemize}[leftmargin=1.2em]
  \item \textbf{Linear Regression $t$-Test (\texttt{LinRegTTest}):}
  1. Enter explanatory data into \texttt{L1} and response data into \texttt{L2}.\\
  2. Press \texttt{[STAT]} $\to$ Arrow to \texttt{TESTS} $\to$ Select \texttt{F:LinRegTTest}.\\
  3. Enter \texttt{Xlist: L1}, \texttt{Ylist: L2}, \texttt{Freq: 1}.\\
  4. Select the alternative hypothesis for $\beta$ and $\rho$ ($\neq 0, < 0,$ or $> 0$).\\
  5. Select \texttt{Calculate}. Displays $t$-statistic, $p$-value, $df = n - 2, a, b, s, r^2,$ and $r$.
  \item \textbf{Linear Regression $t$-Interval (\texttt{LinRegTInt}):}
  1. Press \texttt{[STAT]} $\to$ \texttt{TESTS} $\to$ Select \texttt{G:LinRegTInt}.\\
  2. Enter \texttt{Xlist: L1}, \texttt{Ylist: L2}, and desired \texttt{C-Level} (e.g., 0.95).\\
  3. Select \texttt{Calculate}. Displays the confidence interval for $\beta$, $b_1, df, s, \text{SE}(b_1), a, r^2,$ and $r$.
\end{itemize}
\end{calculatorbox}

\section{Comprehensive Practice Problem Set}

\subsection*{Part I: 20 Multiple-Choice Questions with Explanations}

\begin{enumerate}[leftmargin=1.5em]
  \item What are the degrees of freedom for a linear regression $t$-test with a sample size of $n = 25$?
  \begin{enumerate}[label=(\Alph*)]
    \item 25 \item 24 \item 23 \item 2 \item 1
  \end{enumerate}
  \textbf{Answer: (C).} Regression inference for slope uses $df = n - 2 = 25 - 2 = 23$.

  \item The null hypothesis in a regression significance test evaluating linear association is:
  \begin{enumerate}[label=(\Alph*)]
    \item $H_0: b_1 = 0$ \item $H_0: \beta = 0$ \item $H_0: \mu = 0$ \item $H_0: r = 1$ \item $H_0: \alpha = \beta$
  \end{enumerate}
  \textbf{Answer: (B).} Hypotheses must be stated in terms of the population parameter $\beta$, never the sample statistic $b_1$.

  \item Which of the following conditions ensures that the standard deviation of residuals is constant across all $x$?
  \begin{enumerate}[label=(\Alph*)]
    \item Linear \item Independent \item Normal residuals \item Equal variance \item Random
  \end{enumerate}
  \textbf{Answer: (D).} The Equal Variance condition (homoscedasticity) specifies constant vertical spread of $y$ about the line.

  \item In a regression computer output, the standard error of the slope $\text{SE}(b_1)$ appears under which column header?
  \begin{enumerate}[label=(\Alph*)]
    \item Coef \item SE Coef \item T \item P \item R-sq
  \end{enumerate}
  \textbf{Answer: (B).} The standard error of the slope is located in the ``SE Coef'' column on the row for the explanatory variable.

  \item A 95\% confidence interval for the slope $\beta$ is $(0.45, 1.85)$. Based on this interval:
  \begin{enumerate}[label=(\Alph*)]
    \item There is no linear relationship between $x$ and $y$. \item We fail to reject $H_0: \beta = 0$ at $\alpha = 0.05$. \item We reject $H_0: \beta = 0$ at $\alpha = 0.05$ because zero is not contained in the interval. \item The correlation coefficient must be negative. \item $r^2 > 0.95$.
  \end{enumerate}
  \textbf{Answer: (C).} Because the entire interval is strictly positive and excludes 0, zero is not a plausible value for $\beta$.

  \item If the residual plot exhibits a distinct fan shape (narrow on the left, wide on the right), which condition is violated?
  \begin{enumerate}[label=(\Alph*)]
    \item Linearity \item Independence \item Normality \item Equal variance \item Randomness
  \end{enumerate}
  \textbf{Answer: (D).} A fan or funnel shape violates the Equal Variance assumption.

  \item What does the statistic $s$ represent in a regression output?
  \begin{enumerate}[label=(\Alph*)]
    \item The slope of the line \item The correlation coefficient \item The estimated standard deviation of the residuals \item The sample size \item The standard error of the intercept
  \end{enumerate}
  \textbf{Answer: (C).} $s = \sqrt{\frac{\sum e_i^2}{n-2}}$ estimates the typical distance that observed $y$-values deviate from the predicted line.

  \item A regression yields $b_1 = 3.60$ with $\text{SE}(b_1) = 0.90$ for $n = 20$. The $t$-statistic for $H_0: \beta = 0$ is:
  \begin{enumerate}[label=(\Alph*)]
    \item 0.25 \item 4.00 \item 2.70 \item 3.24 \item 18.0
  \end{enumerate}
  \textbf{Answer: (B).} $t = \frac{b_1 - 0}{\text{SE}(b_1)} = \frac{3.60}{0.90} = 4.00$.

  \item If computer output reports $R\text{-sq} = 64\%$ and the slope coefficient is negative, what is the correlation coefficient $r$?
  \begin{enumerate}[label=(\Alph*)]
    \item 0.64 \item 0.80 \item $-0.80$ \item $-0.64$ \item $-0.4096$
  \end{enumerate}
  \textbf{Answer: (C).} $r = -\sqrt{0.64} = -0.80$. The sign of $r$ matches the sign of the slope $b_1$.

  \item How is the Normality of residuals condition checked in practice?
  \begin{enumerate}[label=(\Alph*)]
    \item By checking whether $n \ge 30$ \item By inspecting a histogram, boxplot, or normal probability plot of the residuals \item By testing $r^2$ \item By plotting $y$ versus $x$ \item By checking $n \le 0.10N$
  \end{enumerate}
  \textbf{Answer: (B).} Normality applies to the distribution of residuals, checked via an NPP or histogram of residuals.

  \item If the $p$-value for the slope $t$-test is $0.003$ and $\alpha = 0.01$, the correct conclusion is:
  \begin{enumerate}[label=(\Alph*)]
    \item Fail to reject $H_0$ \item Reject $H_0$; there is convincing evidence of a linear association \item Accept $H_0$ \item Conclude the relationship is non-linear \item Extrapolate predictions
  \end{enumerate}
  \textbf{Answer: (B).} Since $p\text{-value} = 0.003 < 0.01$, we reject $H_0$.

  \item When constructing a 90\% confidence interval for slope with $n = 12$, what critical value $t^*$ is used?
  \begin{enumerate}[label=(\Alph*)]
    \item $t^*$ with $df = 12$ \item $t^*$ with $df = 11$ \item $t^*$ with $df = 10$ \item $z^* = 1.645$ \item $z^* = 1.960$
  \end{enumerate}
  \textbf{Answer: (C).} $df = n - 2 = 12 - 2 = 10$.

  \item The point of averages $(\bar{x}, \bar{y})$ in simple linear regression:
  \begin{enumerate}[label=(\Alph*)]
    \item Is always an outlier \item Always lies exactly on the least-squares regression line \item Has residual equal to 1 \item Has leverage equal to 1 \item Is never on the line
  \end{enumerate}
  \textbf{Answer: (B).} The LSRL always passes through $(\bar{x}, \bar{y})$ by mathematical definition of the intercept $b_0 = \bar{y} - b_1 \bar{x}$.

  \item In regression, which value measures the percentage of variation in $y$ explained by the linear model?
  \begin{enumerate}[label=(\Alph*)]
    \item $r$ \item $r^2$ \item $s$ \item $\text{SE}(b_1)$ \item $t$
  \end{enumerate}
  \textbf{Answer: (B).} The coefficient of determination $r^2$ measures explained variation.

  \item A two-sided $p$-value in a computer printout is $0.040$. For $H_a: \beta > 0$ with positive $b_1$, the one-sided $p$-value is:
  \begin{enumerate}[label=(\Alph*)]
    \item 0.040 \item 0.080 \item 0.020 \item 0.960 \item 0.010
  \end{enumerate}
  \textbf{Answer: (C).} A one-sided $p$-value in the direction of the sample slope is half the two-sided value: $0.040 / 2 = 0.020$.

  \item If $n = 10$, what are the degrees of freedom for the slope $t$-statistic?
  \begin{enumerate}[label=(\Alph*)]
    \item 8 \item 9 \item 10 \item 18 \item 20
  \end{enumerate}
  \textbf{Answer: (A).} $df = n - 2 = 10 - 2 = 8$.

  \item The sum of the residuals $\sum (y_i - \hat{y}_i)$ for any least-squares regression line is:
  \begin{enumerate}[label=(\Alph*)]
    \item Always positive \item Always negative \item Always exactly zero \item Equal to $r^2$ \item Equal to $s$
  \end{enumerate}
  \textbf{Answer: (C).} A mathematical property of the LSRL is that positive and negative residuals sum to zero.

  \item What does a residual plot with a curved U-shape indicate?
  \begin{enumerate}[label=(\Alph*)]
    \item A linear model is inappropriate \item The correlation is 1 \item The slope is zero \item Conditions are satisfied \item Extrapolation was performed
  \end{enumerate}
  \textbf{Answer: (A).} Any systematic curve in the residual plot indicates that the relationship is non-linear.

  \item If the confidence interval for slope $\beta$ is $(-1.5, 0.5)$, what can be concluded at the corresponding $\alpha$?
  \begin{enumerate}[label=(\Alph*)]
    \item $\beta$ is significantly positive \item $\beta$ is significantly negative \item Fail to reject $H_0: \beta = 0$ because 0 is inside the interval \item $r = 0$ \item Sample size was too large
  \end{enumerate}
  \textbf{Answer: (C).} Since 0 is contained within the confidence interval, zero is a plausible value for $\beta$.

  \item Which acronym summarizes the regression inference conditions?
  \begin{enumerate}[label=(\Alph*)]
    \item BINS \item BITS \item CUSS \item PANIC \item LINER
  \end{enumerate}
  \textbf{Answer: (E).} LINER: Linear, Independent, Normal residuals, Equal variance, Random.
\end{enumerate}

\subsection*{Part II: Score-4 Free-Response Problem with Complete Rubric}

\begin{frqrubricbox}[FRQ 9.1: Comprehensive Regression Inference on Computer Output]
\textbf{Problem:} A sports physiologist studies the relationship between weekly training volume ($x$, in miles) and marathon completion time ($y$, in minutes) for an SRS of $n = 20$ competitive amateur runners. The regression analysis produces:
\begin{center}
\begin{tabular}{lcccc}
\toprule
\textbf{Predictor} & \textbf{Coef} & \textbf{SE Coef} & \textbf{T} & \textbf{P} \\
\midrule
Constant & 310.50 & 14.25 & 21.79 & 0.000 \\
Miles & $-1.85$ & 0.35 & $-5.29$ & 0.000 \\
\bottomrule
\end{tabular}
\end{center}
$S = 8.42\text{ min}, \quad R\text{-sq} = 60.8\%$.\\
(a) State the equation of the least-squares regression line.\\
(b) Interpret the slope $-1.85$ in context.\\
(c) Construct and interpret a 95\% confidence interval for the true slope $\beta$.\\
(d) Explain what $S = 8.42$ minutes means in the context of this study.

\textbf{Grader-Approved Score-4 Solution:}\\
\textbf{Part (a):}\\
$\widehat{\text{Marathon Time}} = 310.50 - 1.85(\text{Weekly Miles})$.\\
\textbf{Part (b):}\\
For each additional 1 mile of weekly training volume, the predicted marathon completion time decreases by approximately 1.85 minutes.\\
\textbf{Part (c):}\\
\textit{Step 1: Parameter \& Name:} We construct a 95\% confidence interval for the true slope $\beta$ of the population regression line relating weekly training miles to marathon finish time. Procedure: Linear regression $t$-interval for slope.\\
\textit{Step 2: Conditions:} Stated as an SRS of 20 runners ($n = 20 < 0.10(\text{All runners})$). Assume scatterplot and residual plot verify linearity, normal residuals, and equal variance.\\
\textit{Step 3: Interval Calculations:} $df = n - 2 = 20 - 2 = 18$. For 95\% confidence with $df = 18$, $t^* = 2.101$.\\
From output: $b_1 = -1.85, \text{SE}(b_1) = 0.35$.\\
\[ \text{CI} = -1.85 \pm 2.101(0.35) = -1.85 \pm 0.7354 = \mathbf{(-2.5854, \; -1.1146)\text{ min/mile}} \]
\textit{Step 4: Interpretation:} We are 95\% confident that for each additional 1 mile per week of training, the true mean marathon completion time decreases by between 1.11 minutes and 2.59 minutes.\\
\textbf{Part (d):}\\
When using weekly training volume to predict marathon completion time with this linear model, predictions will typically deviate from actual observed race times by approximately 8.42 minutes.
\end{frqrubricbox}

\section{Unit 9 Master Summary}
\begin{lessonbox}[Unit 9 Essential Review Checklist]
\begin{itemize}[leftmargin=1.2em]
  \item \textbf{LINER:} Linear, Independent, Normal residuals, Equal variance, Random.
  \item \textbf{Degrees of Freedom:} $df = n - 2$ for all slope $t$-intervals and $t$-tests.
  \item \textbf{Slope CI Formula:} $b_1 \pm t^* \cdot \text{SE}(b_1)$.
  \item \textbf{Slope Test Statistic:} $t = (b_1 - 0)/\text{SE}(b_1)$.
  \item \textbf{Computer Output:} Extract $b_1$ and $\text{SE}(b_1)$ from second row; remember $r$ has the same sign as $b_1$.
\end{itemize}
\end{lessonbox}
\clearpage
